%% prezentacia.tex
%% Copyright 2016, 2017 Juraj Szasz <juraj.szasz3@gmail.com>
%
% This work may be distributed and/or modified under the conditions of the LaTeX
% Project Public License, either version 1.3 of this license or (at your option)
% any later version.  The latest version of this license is in
%   http://www.latex-project.org/lppl.txt
% and version 1.3 or later is part of all distributions of LaTeX version
% 2005/12/01 or later.
%
% This work has the LPPL maintenance status `author-maintained'.
%
% This work consists of the following files: diplomovka.tex, 01-uvod.tex,
% 02-symboly.tex, 03-kapitola1.tex, 04-kapitola2.tex, 05-kapitola3.tex,
% 06-zaver.tex, 07-priloha.tex, literatura.bib, opakuj.sty, ucmthesis.sty and
% ucmplainnat.bst.
%
% The file `ucmplainnat.bst' was derived from the `csplainnat.bst' BibTeX style
% for Czech references style (According to CSN ISO 690) available at:
%   https://github.com/mudrd8mz/csplainnat
%
% All the other files above were derived from the Masaryk University, Faculty of
% Science template available for download at:
%   http://www.sci.muni.cz/NW/STUD/SablonyPraci/Sablona_Pruvodce_(v1.9).zip
%
%%%%%%%%%%%%%%%%%%%%%%%%%%%%%%%%%%%%%%%%%%%%%%%%%%%%%%%%%%%%%%%%%%%%%%%%%%%%%%%%
%%
%%  Sablona Bc+Mgr+RNDr pre FPV UCM
%%  Autor: Juraj Szasz (juraj.szasz3@gmail.com)
%%
%%  Zalozene na sablone PriF MU v Brne
%%    Autor: Petr Zemanek (zemanekp@math.muni.cz)
%%
%%  Typeset in LaTeX-2e
%%
%%%%%%%%%%%%%%%%%%%%%%%%%%%%%%%%%%%%%%%%%%%%%%%%%%%%%%%%%%%%%%%%%%%%%%%%%%%%%%%%

%\documentclass[12pt,a4paper,oneside,final]{book}
\documentclass{beamer}
\usetheme{Warsaw} % pouzitie temy Warsaw

\usepackage[utf8]{inputenc}
\usepackage[LGR,IL2]{fontenc}
\usepackage[english,slovak]{babel} % posledny jazyk v zozname je implicitny

\usepackage{mathptmx}  % volne dostupny font Adobe Times Roman

\usepackage{longtable} % balik pre tabulky presahujuce jednu stranu
                       % (potrebne pre zoznam znaceni)

\usepackage{booktabs}  % balík pre vysadzanie profesionalnych tabuliek

% balik pre postranne popisky obrazkov a tabuliek
\usepackage[leftcaption,ragged]{sidecap}

\usepackage{amsmath,amssymb,amsthm} % baliky pre matematicke vzorce
\usepackage{natbib}    % balik pre harvardsky styl citacii

\usepackage{gfsporson} % balik pre vysadzanie greckeho textu
\usepackage{fancyvrb}  % balik pre oramovanie verbatim prostredia
\usepackage{tabularx}
\usepackage{tikz}

\usebackgroundtemplate{%
  \tikz\node[opacity=0.05]{\parbox{\paperwidth}{\centering\includegraphics[height=\paperheight]{logo/UCM_(color).png}}};
}


%%%%%%%%%%%%%%%%%%%%%%%%%%%%%%%%%%%%%%%%%%%%%%%%%%%%%%%%%%%%%%%%%%%%%%%%%%%%%%%%
%%
%%  Nastavenie stylu citacii
%%
%%%%%%%%%%%%%%%%%%%%%%%%%%%%%%%%%%%%%%%%%%%%%%%%%%%%%%%%%%%%%%%%%%%%%%%%%%%%%%%%
\bibliographystyle{ucmplainnat}
\setcitestyle{authoryear}

%%%%%%%%%%%%%%%%%%%%%%%%%%%%%%%%%%%%%%%%%%%%%%%%%%%%%%%%%%%%%%%%%%%%%%%%%%%%%%%%
%%
%%  Nastavenie sablony
%%
%%%%%%%%%%%%%%%%%%%%%%%%%%%%%%%%%%%%%%%%%%%%%%%%%%%%%%%%%%%%%%%%%%%%%%%%%%%%%%%%

\graphicspath{{obr/}}                 % implicitna cesta pre umiestnenie
                                      % obrazkov v adresari 'obr'

\author{Bc. Juraj Szász}

\institute{\begin{tabular}{l@{\hspace{0.5em}}l}
  \textbf{Študijný program:} & Aplikovaná biológia \\
  \textbf{Študijný obor:}  & Biológia \\
  \textbf{Školiace pracovisko:} & Katedra biológie \\
  \textbf{Školiteľ:} & prof. RNDr. Ján Kraic, PhD \\
\end{tabular}
 \\ \vspace{10pt} \textcolor{blue}{Fakulta prírodných vied} \\ \textcolor{blue}{\textbf{Univerzita sv.\,Cyrila a Metoda v Trnave}}}


\title{Scientometrická analýza FPV~UCM v~Trnave}

\subtitle{Diplomová práca}

\date{25.\,januára 2017}

\setbeamercovered{transparent}


%\NazovKatedry{Katedra biológie}{Department of Biology}

%\MiestoRieseniaPrace{Trnava}{Trnava}

%\AkademickyRok{2016/2017}

%\SkolitelSTitulmi{prof. RNDr. Ján Kraic, PhD}

%\StudijnyProgram{Aplikovaná biológia}{Aplicated Biology}

%\StudijnyObor{Biológia}{Biology}


%\TextAbstraktov{Scientometrická analýza slúži na kvalitatívne a~kvantitatívne
% hodnotenie vedeckej práce.  Vykonáva sa pomocou štatistickej analýzy
% publikačnej činnosti jednotlivých pracovníkov, inštitúcií a~dokonca celých
% krajín.  Kvalitu vedeckých publikácií určuje množstvo citácií (odkazov) na
% publikácie.  Viac citácií na článok vyjadruje jeho popularitu, kvalitu
% a~prínos pre vedecký pokrok.  Cieľom tejto práce je kvalitatítne
% a~kvantitatívne zhodnotiť publikačnú činnosť pracovníkov Fakulty prírodných
% vied Univerzity sv.\,Cyrila a Metoda v~Trnave za~obdobie 2000--2016.}{The
% purpose of scientometric analysis is to provide qualitative and quantitative
% evaluation of the scientific work.  It is being performed by statistical
% analysis of the publishing activity---addressing individual researchers,
% institutions or even whole countries.  The quality of a scientific publication
% is derived from the number of citations.  The amount of references to a paper
% determines its popularity, quality and its impact on scientific progress.  The
% goal of this thesis is to carry out qualitative and quantitative analysis of
% the publishing activity of the academic staff members affiliated to Faculty of
% Natural Science, University of Ss.\,Cyril and Methodeus in Trnava during
% 2000--2016.}

%%%%%%%%%%%%%%%%%%%%%%%%%%%%%%%%%%%%%%%%%%%%%%%%%%%%%%%%%%%%%%%%%%%%%%%%%%%%%%%%
%%
%% Definícia vlastných príkazov
%%
%%%%%%%%%%%%%%%%%%%%%%%%%%%%%%%%%%%%%%%%%%%%%%%%%%%%%%%%%%%%%%%%%%%%%%%%%%%%%%%%

% \newcommand{\Cbb}{\mathbb{C}}
% \newcommand{\Rbb}{\mathbb{R}}
% \newcommand{\Zbb}{\mathbb{Z}}
% \newcommand{\Nbb}{\mathbb{N}}
\newcommand{\tm}{\texttrademark}
\newcommand{\R}{\textsuperscript{\textregistered}}

% Environment for source code
\newenvironment{source}[1][\scriptsize\bfseries]{
  \begin{list}{}{
      \setlength{\leftmargin}{1em}}%
      \setlength{\itemsep}{0pt}%
      \setlength{\parskip}{0pt}%
      \setlength{\parsep}{0pt}%
      \renewcommand{\baselinestretch}{1.0}%
  \item#1}
  {\end{list}}
%%%%%%%%%%%%%%%%%%%%%%%%%%%%%%%%%%%%%%%%%%%%%%%%%%%%%%%%%%%%%%%%%%%%%%%%%%%%%%%%
%%
%%  Dodatocne nastavenia
%%
%%%%%%%%%%%%%%%%%%%%%%%%%%%%%%%%%%%%%%%%%%%%%%%%%%%%%%%%%%%%%%%%%%%%%%%%%%%%%%%%

% Vytvori pomocny subor pre register
%\makeindex

% Zapnutie opakovania matematickych symbolov
% %Podle Tesaříkova czech.ldf
\makeatletter
\def\Deleni{%
 \ifx\protect\@typeset@protect
   \ifhmode
     \ifinner
       \bbl@afterelse\bbl@afterelse\bbl@afterelse\cs@hyphen
     \else
       \bbl@afterfi\bbl@afterelse\bbl@afterelse\cs@firsthyphen
     \fi
   \else
     \bbl@afterfi\bbl@afterelse\cs@hyphen
   \fi
 \else
   \bbl@afterfi\cs@hyphen
 \fi }
\makeatother

%Opakování symbolů binárních operací a relací při zalomení řádku
%Autor: Josef Tkadlec tkadlec@fel.cvut.cz

\relpenalty     =10000      % aby se nelámalo v jiných než ošetřených
\binoppenalty   =10000
\exhyphenpenalty=1000       % aby spíše nouzově (implicitně je 50)
                           % "lokálně" lze zakázat {...}

\def\neq {\mathrel{\not=}}  % aby nedocházelo k lámání \not=/=
\let\ne=\neq

\def\OpakujPrikaz #1#2{\let #2=#1
 \def #1{#2\nobreak\discretionary{}{\hbox{$#2$}}{}}}
\def\OpakujZnak #1#2{\mathchardef #2=\mathcode`#1
 \activedef #1{#2\nobreak\discretionary{}{\hbox{$#2$}}{}}
 \uccode`\~=0 \mathcode`#1="8000 }
%Doplnil Kuben pro nový czech.ldf \expandafter možná nemusí být
\def\OpakujZnakMinus #1#2{\mathchardef #2=\mathcode`#1
 \activedef #1{\ifmmode#2\nobreak\discretionary{}{\hbox{$#2$}}{}\else\expandafter\Deleni\fi }
 \uccode`\~=0 \mathcode`#1="8000 }
\def\activedef #1{\uccode`\~=`#1 \uppercase{\def~}}

\OpakujPrikaz {\neq }{\neqORI}  \let \ne=\neq
\OpakujPrikaz {\leq }{\leqORI}  \let \le=\leq
\OpakujPrikaz {\geq }{\geqORI}  \let \ge=\geq
\OpakujPrikaz {\cup }{\cupORI}
\OpakujPrikaz {\cap }{\capORI}
\OpakujPrikaz {\times }{\timesORI}
\OpakujPrikaz {\subset }{\subsetORI}
\OpakujPrikaz {\subseteq }{\subseteqORI}
\OpakujPrikaz {\supset }{\supsetORI}
\OpakujPrikaz {\supseteq }{\supseteqORI}

\OpakujPrikaz {\cdot }{\cdotORI}
\OpakujPrikaz {\setminus }{\setminusORI}

\OpakujZnak <{\lessORI}
\OpakujZnak >{\greaterORI}
\OpakujZnak +{\plusORI}
\AtBeginDocument {\OpakujZnak ={\eqORI} \OpakujZnakMinus -{\minusORI}}

%%%%%%%%%%%%%%%%%%%%%%%%%%%%%%%%%%%%%%%%%%%%%%%%%%%%%%%%%%%%%%%%%%%%%%%%%%%%%%%%
%%
%%  Zaciatok dokumentu
%%
%%%%%%%%%%%%%%%%%%%%%%%%%%%%%%%%%%%%%%%%%%%%%%%%%%%%%%%%%%%%%%%%%%%%%%%%%%%%%%%%

% [0] Názov školy, fakulty, práce... to isté ako na prvej strany diplomovky
% [1] Obsah
% [2] Čo je to scientometria? nejaky funny obrázok.
% [2.1] Čo je to scientometria? definícia
% [2.2] Čo je to scientometria? obrázok rozdelenia informetrie
% [3.1] Scopus: fakty o veľkosti a rozdelení
% [3.2] Scopus: screenshot
% [4.1] Web of Science: fakty o veľkosti a rozdelení
% [4.2] Web of Science: screenshot
% [5.1] Google Scholar: nejaké fakty
% [5.2] Google Scholar: screenshot
% [5.3] Google Scholar: príklad chýb
% [6] Citačné indikátory tabuľka
% [7.0] Praktická časť: Datamining nejaký funny obrazok
% [7.1.1] Scopus: dopyt
% [7.1.2] Scopus: screenshot
% [7.2.1] Web of Science: dopyt
% [7.2.2] Web of Science: screenshot
% [7.3.1] Analýza pracovníkov: zoznam pracovníkov
% [7.3.2] Analýza pracovníkov chýbajúce dáta
% [7.4.1] Program Publish or Perish: nejaké info
% [7.4.2] Program Publish or Perish: screenshot
% [7.5.1] Skripty na spracovanie dát: nejaké info
% [8] Výsledky: nejaký funny obrázok
% [8.1] Výsledky: Vývoj publikačnej činnosti FPV
% [8.2] Výsledky: Vývoj citovanisti publikácií FPV
% [8.3] Výsledky: 1. strana tabuľky publikácií
% [8.4.1]--[8.4.16] Výsledky grafy v poradí, akom sú v diplomovke
% [9] Záver posledný odstavec práce
% [10] Ďakujem za pozornosť. Diskusia a nejaký funny obrázok (napr. PhD comics)
% [11] Použitá literatúra

\begin{document}

%%
%%  Uvodna strana
%%

\frame{\maketitle}

\frame{\tableofcontents}


\section[Úvod]{Úvod}

\subsection{Definícia scientometrie}


\begin{frame}
\frametitle{Čo je scientometria?}

\begin{block}{Definícia scientometrie}
Scientometria je....
\end{block}

\end{frame}

\subsection{Citačné registre (Scopus, Web of Science, Google Scholar)}

\begin{frame}
\frametitle{Scopus}

Contents\dots

\end{frame}

\begin{frame}
\frametitle{Web of Science}

Contents\dots

\end{frame}

\begin{frame}
\frametitle{Google Scholar}


Contents\dots


\end{frame}


\subsection{Citačné indikátory}

\begin{frame}
\frametitle{Prehľad citačných indikátorov}


Contents\dots


\end{frame}

\section{Ciele práce}

\begin{frame}
\frametitle{Ciele práce}

\begin{enumerate}
\item<1-> Získať bibliografické záznamy všetkých publikácií každého vedeckého
  pracovníka, ktorý je uvedený na stránkach fakulty FPV UCM v Trnave a na
  stránkach príslušných katedier.
\item<2-> Získať bibliografické záznamy pre práce, v ktorých má aspoň jeden
  spoluautor príslušnosť FPV UCM v Trnave.
\item<3-> Urobiť citačnú analýzu vedeckých publikácií pracovníkov FPV UCM v
  Trnave.
\item<4-> Rozdeliť citačné analýzy jednotlivých pracovníkov podľa príslušnosti
  ku katedrám a vytvoriť sumárne citačné analýzy pre každú katedru.
\item<5-> Graficky zhodnotiť a porovnať výsledok citačných analýz pre každu
  katedru a graficky zhodnotiť vývoj publikačnej činnosti FPV UCM v Trnave.
\end{enumerate}

\end{frame}


\section{Materiály a metódy}

\subsection{Získavanie dát}

\begin{frame}
\frametitle{Scopus}

Contents\dots


\end{frame}

\begin{frame}
\frametitle{Web of Science}

Contents\dots


\end{frame}

\subsection{Analýza pracovníkov}

\begin{frame}
\frametitle{Zoznam pracovníkov}

Contents\dots


\end{frame}

\begin{frame}
\frametitle{Chýbajúce dáta}

Contents\dots


\end{frame}


\subsection{Program Publish or Perish}

\begin{frame}
\frametitle{Program Publish or Perish}


Contents\dots


\end{frame}


\subsection{Vlastné programy na spracovanie dát}

\begin{frame}
\frametitle{Vlastné programy na spracovanie dát}

\begin{block}{scientometry-data-proc}

Contents\dots

\end{block}

\begin{block}{scientometry-plot-gen}

Contents\dots

\end{block}

\end{frame}


\section{Výsledky a diskusia}

\subsection{Vývoj publikačnej činnosti FPV}


\begin{frame}


Contents\dots


\end{frame}


\subsection{Vývoj citovanosti publikácii FPV}


\begin{frame}


Contents\dots


\end{frame}


\subsection{Tabuľka publikácii}


\begin{frame}


Contents\dots


\end{frame}


\section{Záver}

\begin{frame}
\frametitle{Záver}

\begin{itemize}
\item<1-> Z technických dôvodov sme do analýz nezahrnuli odstránenie
  autocitácií.
\item<2-> Práca nezahŕňa výpočet spoluatorských sietí, pretože by tak presiahla
  predpísaný rozsah.
\item<3-> I napriek sa nám v rámci možností podarilo splniť stanovené ciele.
\item<4-> Táto práca tak predstavuje prvú komplexnejšiu scientometrickú analýzu
  FPV UCM v Trnave a môže sa stať užitočným podkladom pre ďaľšie práce tohto
  typu.
\end{itemize}

\end{frame}


\begin{frame}[plain]
  \begin{center}
    \Huge Ďakujem za pozornosť!
    \vfill
    \Large Juraj Szász
  \end{center}
\end{frame}

\begin{frame}[plain]
  \frametitle{Iné zdroje impakt faktoru}
  \begin{table}
    \small\centering
    \begin{tabular}{lccl}
      \toprule
      Časopis & ISI JIF & IF & Zdroj \\
      \midrule
      Polish Journal of Env. Studies  & 0,79  & 0,71           & Research Gate$^{\mathrm{a}}$ \\[0.5ex]
      Pharmazie                       & 1,264 & 1,32           & Research Gate$^{\mathrm{b}}$ \\[0.5ex]
      Applied Mechanics and Materials & --    & 0,16           & Research Gate$^{\mathrm{c}}$ \\[0.5ex]
      Advanced Materials Research     & --    & 0,23           & Research Gate$^{\mathrm{d}}$ \\[0.5ex]
      Nova Biotechnologica et Chimica & --    & 0,188$^\dagger$ & De Gruyter$^{\mathrm{e}}$    \\[0.5ex]
      Comm. in Comp. and Inf. Science & --    & 0,35           & Research Gate$^{\mathrm{f}}$ \\[0.5ex]
      \bottomrule
      \multicolumn{4}{l}{\tiny $^\dagger$ Hodnota nie je impakt faktor, ale impakt na publikáciu (IPP)} \\
      \multicolumn{4}{l}{\tiny $^{\mathrm{a}}$ \url{https://www.researchgate.net/journal/1230-1485\_Polish\_Journal\_of\_Environmental\_Studies}} \\
      \multicolumn{4}{l}{\tiny $^{\mathrm{b}}$ \url{https://www.researchgate.net/journal/0031-7144\_Pharmazie}} \\
      \multicolumn{4}{l}{\tiny $^{\mathrm{c}}$ \url{https://www.researchgate.net/journal/1660-9336\_Applied\_Mechanics\_and\_Materials}} \\
      \multicolumn{4}{l}{\tiny $^{\mathrm{d}}$ \url{https://www.researchgate.net/journal/1662-8958\_Advanced\_Materials\_Research}} \\
      \multicolumn{4}{l}{\tiny $^{\mathrm{e}}$ \url{https://www.degruyter.com/view/j/nbec}} \\
      \multicolumn{4}{l}{\tiny $^{\mathrm{f}}$ \url{https://www.researchgate.net/journal/1865-0929\_Communications\_in\_Computer\_and\_Information\_Science}} \\
    \end{tabular}
  \end{table}
\end{frame}

%%
%%  Zoznam pouzitej literatury
%%
% \begin{frame}
%   \frametitle{Zoznam literatúry}
% \bibliography{literatura}{}
% \end{frame}

\end{document}
